\documentclass[aspectratio=43]{beamer}

% ==============================================================================
% TEMA Y ESTILO ACADÉMICO MODERNO
% ==============================================================================
\usetheme{Madrid}
\usecolortheme{default}

% Paleta de colores personalizada (Azul Marino / Blanco / Gris Técnico)
\definecolor{UBlue}{RGB}{14, 43, 92}
\definecolor{UBlueLight}{RGB}{30, 75, 145}
\definecolor{UAccent}{RGB}{210, 80, 20}
\definecolor{CodeBG}{RGB}{245, 246, 248}
\definecolor{BorderColor}{RGB}{200, 205, 215}

\setbeamercolor{palette primary}{bg=UBlue,fg=white}
\setbeamercolor{palette secondary}{bg=UBlueLight,fg=white}
\setbeamercolor{palette tertiary}{bg=UBlue,fg=white}
\setbeamercolor{palette quaternary}{bg=UBlue,fg=white}
\setbeamercolor{structure}{fg=UBlue}
\setbeamercolor{frametitle}{bg=UBlue,fg=white}
\setbeamercolor{block title}{bg=UBlueLight,fg=white}
\setbeamercolor{block body}{bg=CodeBG,fg=black}
\setbeamercolor{alertblock title}{bg=UAccent,fg=white}
\setbeamercolor{alertblock body}{bg=orange!10,fg=black}

% Configuración de tipografía y elementos
\setbeamertemplate{navigation symbols}{}
\setbeamertemplate{itemize items}[circle]
\setbeamertemplate{enumerate items}[default]

% ==============================================================================
% PAQUETES
% ==============================================================================
\usepackage[utf8]{inputenc}
\usepackage[spanish]{babel}
\usepackage{amsmath,amssymb}
\usepackage{booktabs}
\usepackage{listings}
\usepackage{tikz}
\usetikzlibrary{shapes.geometric, arrows.meta, positioning, fit, backgrounds, calc}

% Configuración de entorno de código (Listings)
\lstset{
    backgroundcolor=\color{CodeBG},
    basicstyle=\ttfamily\fontsize{6.8pt}{7.8pt}\selectfont,
    breakatwhitespace=false,
    breaklines=true,
    captionpos=b,
    frame=single,
    rulecolor=\color{BorderColor},
    keywordstyle=\color{UBlueLight}\bfseries,
    stringstyle=\color{UAccent},
    commentstyle=\color{gray}\itshape,
    showstringspaces=false,
    numbers=none,
    tabsize=2,
    xleftmargin=2pt,
    xrightmargin=2pt
}

% ==============================================================================
% METADATOS DE LA PRESENTACIÓN
% ==============================================================================
\title[Monitoreo Geoespacial de Ganado IoT]{Sistema de Monitoreo Geoespacial y Disuasión de Ganado en Zonas Rurales Sin Cobertura}
\subtitle{Arquitectura Desacoplada de 3 Capas: Edge, Gateway y Servidor Docker}
\author[Trabajo Final]{Defensa de Trabajo Final}
\institute[UNLP]{
    Facultad de Informática\\
    Universidad Nacional de La Plata
}
\date{\today}

% ==============================================================================
% INICIO DEL DOCUMENTO
% ==============================================================================
\begin{document}

% ------------------------------------------------------------------------------
% DIAPOSITIVA 1: PORTADA
% ------------------------------------------------------------------------------
\begin{frame}
    \titlepage
\end{frame}

% ------------------------------------------------------------------------------
% DIAPOSITIVA 2: AGENDA
% ------------------------------------------------------------------------------
\begin{frame}{Agenda de la Presentación}
    \tableofcontents
\end{frame}

% ==============================================================================
% SECCIÓN 1: CONTEXTO Y PROBLEMÁTICA
% ==============================================================================
\section{Contexto y Problemática}

% ------------------------------------------------------------------------------
% DIAPOSITIVA 3: PROBLEMA
% ------------------------------------------------------------------------------
\begin{frame}{Contexto y Problemática en la Ganadería Extensiva}
    \begin{block}{Restricciones Críticas del Entorno Rural}
        \small
        \begin{itemize}
            \item \textbf{Inexistencia de Cobertura Celular:} Las zonas rurales de pastoreo extensivo carecen de conectividad 4G/LTE y redes eléctricas fijas.
            \item \textbf{Autonomía Energética Extrema:} Los dispositivos adosados al animal (collares) deben operar durante meses con baterías de tamaño y peso acotados.
            \item \textbf{Latencia en la Contención:} La detección tradicional por satélite no permite una reacción inmediata cuando el animal traspasa un límite de propiedad.
        \end{itemize}
    \end{block}

    \begin{alertblock}{Objetivo Tecnológico}
        \small
        Diseñar e implementar un sistema que combine \textbf{telemetría geoespacial de ultra bajo consumo}, \textbf{disuasión acústica instantánea en el borde} y \textbf{orquestación centralizada} sobre radiofrecuencia libre (LoRa 433~MHz).
    \end{alertblock}
\end{frame}

% ------------------------------------------------------------------------------
% DIAPOSITIVA 4: SOLUCIÓN ARQUITECTÓNICA
% ------------------------------------------------------------------------------
\begin{frame}{Solución Propuesta: Desacoplamiento en 3 Capas}
    \small
    Para superar estas restricciones, el sistema distribuye las responsabilidades en tres capas especializadas:
    
    \vspace{0.2cm}
    \begin{table}
        \centering
        \scriptsize
        \begin{tabular}{p{2.2cm} p{2.8cm} p{4.6cm}}
            \toprule
            \textbf{Capa} & \textbf{Hardware} & \textbf{Responsabilidad Principal} \\
            \midrule
            \textbf{1. Edge (Collar)} & ESP32 + NEO-6M + SX1278 + Zumbador & Adquisición GPS, Deep Sleep, cálculo local de Bounding Box y disuasión autónoma. \\
            \addlinespace
            \textbf{2. Gateway} & ESP32 + SX1278 + Wi-Fi & Traducción LoRa $\rightarrow$ MQTT, medición física (RSSI/SNR) y auto-descubrimiento. \\
            \addlinespace
            \textbf{3. Servidor Local} & Nginx + Node-RED + InfluxDB + Mosquitto & Ingesta temporal, Ray-Casting ($\mathcal{O}(N)$), mapa Web HUD y bot de Telegram. \\
            \bottomrule
        \end{tabular}
    \end{table}
\end{frame}

% ==============================================================================
% SECCIÓN 2: ARQUITECTURA DEL SISTEMA
% ==============================================================================
\section{Arquitectura del Sistema}

% ------------------------------------------------------------------------------
% DIAPOSITIVA 5: DIAGRAMA GENERAL DE ARQUITECTURA
% ------------------------------------------------------------------------------
\begin{frame}{Arquitectura General del Sistema}
    \begin{figure}
        \centering
        \resizebox{0.96\textwidth}{!}{
        \begin{tikzpicture}[
            node distance=1.4cm and 1.6cm,
            box/.style={draw, rounded corners=2pt, minimum height=1.0cm, align=center, font=\footnotesize, fill=white, line width=0.7pt},
            edgebox/.style={box, draw=UBlue, fill=UBlue!5},
            gwbox/.style={box, draw=UAccent, fill=UAccent!5},
            srvbox/.style={box, draw=black, fill=CodeBG},
            arrow/.style={->, >=Stealth, line width=0.8pt, color=UBlue}
        ]
            % Nodos de Capa Edge
            \node[edgebox] (gps) {GPS\\NEO-6M};
            \node[edgebox, right=0.6cm of gps] (edge) {ESP32 Edge\\(Deep Sleep + RTC)};
            \node[edgebox, right=0.6cm of edge] (sx_tx) {Transceptor\\SX1278 (TX)};
            \node[edgebox, below=0.4cm of edge] (buzzer) {Zumbador\\PWM};

            % Nodos de Capa Gateway
            \node[gwbox, right=1.4cm of sx_tx] (sx_rx) {Transceptor\\SX1278 (RX)};
            \node[gwbox, right=0.6cm of sx_rx] (gw) {ESP32 Gateway\\(mDNS + Wi-Fi)};

            % Nodos de Capa Servidor (Docker)
            \node[srvbox, below=1.1cm of gw] (mqtt) {Mosquitto\\MQTT :1883};
            \node[srvbox, left=1.0cm of mqtt] (nodered) {Node-RED\\:1880};
            \node[srvbox, left=1.0cm of nodered] (influx) {InfluxDB\\:8086};
            \node[srvbox, below=0.7cm of nodered] (nginx) {Nginx Proxy\\:80 (CORS libre)};

            % Nodos Externos
            \node[box, fill=yellow!10, left=1.2cm of nginx] (frontend) {Web HUD\\React / Leaflet};
            \node[box, fill=green!10, right=1.2cm of nginx] (telegram) {Telegram Bot\\Alertas};

            % Conexiones Edge
            \draw[arrow] (gps) -- (edge);
            \draw[arrow] (edge) -- (sx_tx);
            \draw[arrow] (edge) -- (buzzer);

            % Conexión LoRa
            \draw[arrow, color=UAccent, dashed, line width=1pt] (sx_tx) -- node[above, font=\tiny] {LoRa 433~MHz} node[below, font=\tiny] {JSON compact} (sx_rx);

            % Conexiones Gateway
            \draw[arrow] (sx_rx) -- (gw);
            \draw[arrow] (gw) -- node[right, font=\tiny] {Wi-Fi / TCP} (mqtt);

            % Conexiones Servidor
            \draw[arrow] (mqtt) -- node[above, font=\tiny] {geofence/heartbeat} (nodered);
            \draw[arrow] (nodered) -- node[above, font=\tiny] {series temporales} (influx);
            \draw[arrow, <->] (nginx) -- (nodered);
            \draw[arrow, <->] (nginx) -- (influx);

            % Conexiones Externas
            \draw[arrow, <->] (frontend) -- node[above, font=\tiny] {HTTP :80} (nginx);
            \draw[arrow] (nodered) -| (telegram);

            % Contornos de Capas
            \begin{scope}[on background layer]
                \node[draw=UBlue, dashed, fit=(gps)(edge)(sx_tx)(buzzer), label={[font=\scriptsize\bfseries, color=UBlue]above:Capa 1: Módulo Edge (Collar)}] {};
                \node[draw=UAccent, dashed, fit=(sx_rx)(gw), label={[font=\scriptsize\bfseries, color=UAccent]above:Capa 2: Gateway}] {};
                \node[draw=gray, dashed, fit=(mqtt)(nodered)(influx)(nginx), label={[font=\scriptsize\bfseries]above:Capa 3: Servidor en Docker Compose}] {};
            \end{scope}
        \end{tikzpicture}
        }
    \end{figure}
\end{frame}

% ------------------------------------------------------------------------------
% DIAPOSITIVA 6: CAPA EDGE Y GESTIÓN DE ENERGÍA
% ------------------------------------------------------------------------------
\begin{frame}[fragile]{Capa 1: Módulo Edge y Preservación de Contexto}
    \begin{columns}[T]
        \begin{column}{0.48\textwidth}
            \begin{block}{Estrategia de Deep Sleep}
                \scriptsize
                \begin{itemize}
                    \item El microcontrolador apaga CPU y periféricos de radio el 99\% del tiempo.
                    \item Despierta por temporizador RTC cada $N$ minutos o por movimiento.
                    \item El tiempo activo en aire (\textit{Time-on-Air}) se minimiza enviando un payload $<100$ bytes.
                \end{itemize}
            \end{block}
        \end{column}
        \begin{column}{0.48\textwidth}
            \begin{block}{Memoria Estática (\texttt{RTC\_DATA\_ATTR})}
                \scriptsize
                Preserva las coordenadas y límites sin gastar energía en re-inicializaciones:
            \end{block}
            \begin{lstlisting}[language=C++]
// rtc_state.h
extern RTC_DATA_ATTR double rtc_lat;
extern RTC_DATA_ATTR double rtc_lon;
extern RTC_DATA_ATTR uint32_t wake_count;
extern RTC_DATA_ATTR double min_lat;
extern RTC_DATA_ATTR double max_lat;
            \end{lstlisting}
        \end{column}
    \end{columns}
\end{frame}

% ------------------------------------------------------------------------------
% DIAPOSITIVA 7: ALGORITMO GEOMÉTRICO EN EL BORDE
% ------------------------------------------------------------------------------
\begin{frame}[fragile]{Algoritmo Geométrico en el Borde: Bounding Box ($\mathcal{O}(1)$)}
    \begin{block}{¿Por qué un Bounding Box local en el collar?}
        \small
        El collar almacena únicamente los 4 extremos cardinales (\texttt{min\_lat}, \texttt{max\_lat}, \texttt{min\_lon}, \texttt{max\_lon}) del potrero para ahorrar RAM y evitar sincronizaciones pesadas de polígonos por LoRa.
    \end{block}

    \begin{columns}[T]
        \begin{column}{0.54\textwidth}
            \begin{lstlisting}[language=C++]
// Verificacion local O(1)
bool geofence_is_inside(double lat, double lon,
                        double min_lat, double max_lat,
                        double min_lon, double max_lon) {
  if (lat == 0.0 && lon == 0.0) return true;
  return (lat >= min_lat && lat <= max_lat &&
          lon >= min_lon && lon <= max_lon);
}
            \end{lstlisting}
        \end{column}
        \begin{column}{0.42\textwidth}
            \begin{alertblock}{Autonomía Disuasoria}
                \scriptsize
                Si \texttt{geofence\_is\_inside()} retorna \texttt{false}, el collar enciende el \textbf{zumbador PWM} de inmediato, logrando disuasión acústica real \textbf{sin depender de la nube}.
            \end{alertblock}
        \end{column}
    \end{columns}
\end{frame}

% ------------------------------------------------------------------------------
% DIAPOSITIVA 8: CAPA GATEWAY Y AUTO-DESCUBRIMIENTO
% ------------------------------------------------------------------------------
\begin{frame}[fragile]{Capa 2: Módulo Gateway y Auto-descubrimiento}
    \begin{block}{Traducción de Protocolo (LoRa $\rightarrow$ MQTT)}
        \small
        El Gateway recibe las tramas en 433~MHz, mide la calidad de canal (RSSI/SNR) y encapsula los datos hacia Mosquitto en el tópico \texttt{geofence/heartbeat}:
    \end{block}

    \begin{columns}[T]
        \begin{column}{0.5\textwidth}
            \begin{lstlisting}[language=C++]
// gateway-module.cpp: Enriquecimiento
snprintf(mqtt_payload, sizeof(mqtt_payload),
  "{\"device_id\":\"%s\",\"lat\":%.6f,"
  "\"lon\":%.6f,\"fence_status\":\"%s\","
  "\"rssi\":%d,\"snr\":%.2f}",
  dev_id, lat, lon, status, rssi, snr);
mqtt_init_and_publish(mqtt_payload);
            \end{lstlisting}
        \end{column}
        \begin{column}{0.46\textwidth}
            \begin{block}{Auto-descubrimiento (mDNS / Beacon)}
                \scriptsize
                Evita IPs fijas en código ante cambios de router o entorno LAN:
                \begin{itemize}
                    \item Escucha Beacons / UDP Broadcast.
                    \item Header: \texttt{GEOFENCE\_MQTT:<IP>:<PORT>}
                    \item Conexión dinámica al broker.
                \end{itemize}
            \end{block}
        \end{column}
    \end{columns}
\end{frame}

% ------------------------------------------------------------------------------
% DIAPOSITIVA 9: CAPA SERVIDOR DOCKER Y NGINX
% ------------------------------------------------------------------------------
\begin{frame}{Capa 3: Servidor en Docker Compose y Reverse Proxy}
    \small
    Los servicios backend corren en una red bridge privada y se exponen al cliente web mediante un \textbf{Reverse Proxy en Nginx (Puerto 80)}:

    \vspace{0.15cm}
    \begin{table}
        \centering
        \scriptsize
        \begin{tabular}{p{2.6cm} p{3.2cm} p{4.0cm}}
            \toprule
            \textbf{Petición Externa (:80)} & \textbf{Destino Interno (Docker)} & \textbf{Propósito Técnico} \\
            \midrule
            \texttt{GET /} & \texttt{geofence-frontend:80} & Archivos estáticos del mapa React. \\
            \texttt{POST /api/update-fence} & \texttt{geofence-nodered:1880} & Actualizar geometría de geocerca. \\
            \texttt{GET /api/fence} & \texttt{geofence-nodered:1880} & Consultar perímetro en curso. \\
            \texttt{POST /query} & \texttt{geofence-influxdb:8086} & Consultas InfluxQL directas. \\
            \bottomrule
        \end{tabular}
    \end{table}

    \begin{alertblock}{Eliminación de Errores CORS}
        \scriptsize
        Al unificar todo el tráfico bajo el puerto 80 en \texttt{nginx.conf}, el navegador procesa las peticiones como llamadas \textit{Same-Origin}, eliminando cabezales CORS complejos y aislamiento.
    \end{alertblock}
\end{frame}

% ==============================================================================
% SECCIÓN 3: FLUJO DE DATOS Y ALGORITMOS DUALES
% ==============================================================================
\section{Flujo de Datos y Algoritmo Dual}

% ------------------------------------------------------------------------------
% DIAPOSITIVA 10: INGESTIÓN DE SERIES TEMPORALES
% ------------------------------------------------------------------------------
\begin{frame}[fragile]{Flujo de Datos: Ingestión en Node-RED}
    \begin{block}{Tábula \texttt{flow\_tab\_mqtt\_influx} en \texttt{flows.json}}
        \scriptsize
        1. El nodo \texttt{mqtt in} captura el mensaje de \texttt{geofence/heartbeat}.\\
        2. El nodo función \texttt{Parsear Heartbeat} transforma el JSON para persistirlo en InfluxDB:
    \end{block}

    \begin{columns}[T]
        \begin{column}{0.48\textwidth}
            \begin{lstlisting}[language=C++]
// Estructura generada por Node-RED
var fields = {
  latitude: parseFloat(raw.lat),
  longitude: parseFloat(raw.lon),
  rssi: parseInt(raw.rssi),
  snr: parseFloat(raw.snr),
  battery: parseFloat(raw.battery),
  wake_count: parseInt(raw.wake_count)
};
var tags = { device_id: raw.device_id };
            \end{lstlisting}
        \end{column}
        \begin{column}{0.48\textwidth}
            \begin{block}{Persistencia en InfluxDB}
                \scriptsize
                Escribe en la medición \textbf{\texttt{cattle\_position}} (\texttt{geofence\_db}).
                \vspace{0.15cm}
                
                Permite consultar historiales y trayectorias exactas por dispositivo en ventanas móviles (\texttt{now() - 24h}).
            \end{block}
        \end{column}
    \end{columns}
\end{frame}

% ------------------------------------------------------------------------------
% DIAPOSITIVA 11: ALGORITMO EN SERVIDOR Y ALERTAS
% ------------------------------------------------------------------------------
\begin{frame}{Algoritmo en Servidor: Ray-Casting ($\mathcal{O}(N)$) y Alertas}
    \begin{block}{¿Por qué calcular nuevamente la geocerca en el servidor?}
        \small
        Mientras el collar usa un Bounding Box ($\mathcal{O}(1)$) para contención acústica, Node-RED evalúa las coordenadas contra el polígono irregular de $N$ vértices mediante \textbf{Ray-Casting (Point-in-Polygon)} para determinar la \textbf{verdad canónica}.
    \end{block}

    \begin{columns}[T]
        \begin{column}{0.5\textwidth}
            \small
            \begin{itemize}
                \item \textbf{Sondeo Continuo:} Cron en Node-RED consulta InfluxDB cada 15~s.
                \item \textbf{Rate Limiting ($40\text{ s}$):} Suprime el envío repetitivo de alarmas si el animal camina sobre el límite.
            \end{itemize}
        \end{column}
        \begin{column}{0.48\textwidth}
            \begin{alertblock}{Gestión por Telegram Bot}
                \scriptsize
                Si hay violación geométrica verídica y el dispositivo no está silenciado (\texttt{cattle\_mute\_status} / comandos \texttt{/mute}), se emite una alerta REST hacia la API de Telegram con botones interactivos.
            \end{alertblock}
        \end{column}
    \end{columns}
\end{frame}

% ==============================================================================
% SECCIÓN 4: CONCLUSIONES Y DEFENSA
% ==============================================================================
\section{Conclusiones y Defensa}

% ------------------------------------------------------------------------------
% DIAPOSITIVA 12: CONCLUSIONES
% ------------------------------------------------------------------------------
\begin{frame}{Conclusiones y Viabilidad Técnica}
    \begin{block}{Aportes Principales de la Ingeniería}
        \small
        \begin{enumerate}
            \item \textbf{Autonomía Resiliente:} La separación entre Bounding Box ($\mathcal{O}(1)$) en el collar y Ray-Casting ($\mathcal{O}(N)$) en el servidor garantiza contención acústica aún ante pérdida total del enlace LoRa.
            \item \textbf{Eficiencia Energética y de Espectro:} El uso de \texttt{RTC\_DATA\_ATTR} con Deep Sleep y paquetes LoRa compactos minimiza el consumo eléctrico y el \textit{Time-on-Air}.
            \item \textbf{Infraestructura Limpia:} Docker Compose junto a Nginx unifican la seguridad, simplifican la red y eliminan colisiones CORS en la web.
        \end{enumerate}
    \end{block}

    \vspace{0.2cm}
    \begin{center}
        \Large \textbf{¡Muchas Gracias!}\\
        \normalsize Espacio abierto para preguntas del jurado.
    \end{center}
\end{frame}

\end{document}
